This chapter fixes the character-theoretic notation used later.

\section{The Galois group}

\begin{definition}\label{def:galois-cyclotomic}
  \lean{KummerCriterion.cyclotomicGalEquivZMod}
  \lean{KummerCriterion.cyclotomicSigmaOfUnit}
  \lean{KummerCriterion.cyclotomicGalEquivZMod_sigmaOfUnit}
  \lean{KummerCriterion.cyclotomicSigmaOfUnit_apply_zeta}
  \leanok
The extension $\Qzetap/\Q$ is Galois of degree $p-1$. The assignment
\[
  \sigma_a(\zetap)=\zetap^a
  \qquad (a\in (\Z/p\Z)^\times)
\]
identifies
\(
  \Gal{\Qzetap/\Q}
\)
with
\(
  (\Z/p\Z)^\times.
\)
\end{definition}

\section{Dirichlet characters}

\begin{definition}\label{def:dirichlet-character}
  \lean{DirichletCharacter}\leanok
Let $N\ge 1$ and let $R$ be a commutative ring. A \emph{Dirichlet character
modulo $N$} with values in $R$ is a multiplicative map
\(
  \chi : \Z/N\Z \to R
\)
which vanishes on non-units.
\end{definition}

When $p$ is prime, the non-trivial characters modulo $p$ are exactly the
characters of the cyclic group $(\Z/p\Z)^\times$. In particular they form a
cyclic group of order $p-1$, and every non-trivial character satisfies the
orthogonality relation
\[
  \sum_{a\in \Z/p\Z} \chi(a) = 0.
\]

\begin{definition}[Teichm\"uller character]
  \label{def:teichmuller}
  \lean{KummerCriterion.teichmullerCharQp}\leanok
  \uses{def:dirichlet-character}
The \emph{Teichm\"uller character} is the unique character
\[
  \omega : (\Z/p\Z)^\times \longrightarrow \Z_p^\times
\]
such that \(\omega(a)\equiv a\pmod p\) for every nonzero residue class
\(a\). It generates the character group.
\end{definition}

Thus every Dirichlet character modulo $p$ can be written uniquely as a power
\(\omega^j\) with $0\le j\le p-2$.

\section{Even and odd characters}

\begin{definition}\label{def:even-odd}
  \lean{DirichletCharacter.Even}\leanok
  \uses{def:dirichlet-character}
A Dirichlet character $\chi$ is \emph{even} if $\chi(-1)=1$ and \emph{odd}
if $\chi(-1)=-1$.
\end{definition}

Exactly half of the characters modulo $p$ are even and exactly half are odd.
The even characters are precisely those that factor through the quotient
\(
  (\Z/p\Z)^\times / \{\pm 1\},
\)
which is the Galois group of the maximal real subfield $\Qzetapplus$.
This parity distinction is the source of the even/odd decomposition in the
analytic class number formula.

Finally, because $p$ is prime, every non-trivial Dirichlet character modulo
$p$ is primitive. Hence the functional equation for primitive Dirichlet
$L$-functions applies to every odd character that appears later.
